\title{When Do Sustainability Priors Help Reinforcement Learning? A Hybrid FBWM-FTOPSIS-PPO Framework for Dynamic Supplier Order Allocation}

% Add your authors
% Use \author[inst1]{Name} for authors at the same institution
% Use \corref{cor1} for the corresponding author
\author[inst1]{Ali Vaezi\corref{cor1}}
\ead{ali.vaezi@studenti.polito.it}

\author[inst2]{Erfan Rabbani}
\ead{erfaninja@gmail.com}

% Add \cortext
\cortext[cor1]{Corresponding author.}

\affiliation[inst1]{
    organization={Interuniversity Department
of Regional and Urban Studies and Planning, Politecnico di Torino},
    addressline={Corso Duca degli Abruzzi, 24}, 
    city={Turin},
    postcode={10129}, 
    country={Italy}
}
\affiliation[inst2]{
    organization={Department of Management and Production Engineering, Politecnico di Torino},
    addressline={Corso Duca degli Abruzzi, 24}, 
    city={Turin},
    postcode={10129}, 
    country={Italy}
}

% Add your abstract
\begin{abstract}
Global supply chains must balance cost efficiency, sustainability, and resilience under growing demand uncertainty and disruption risk, yet existing approaches address supplier evaluation and dynamic ordering in isolation.
Static multi-criteria decision-making (MCDM) methods rank suppliers effectively but cannot adapt to changing conditions, while reinforcement learning (RL) agents learn adaptive policies but must rediscover supplier quality from scratch.
This paper proposes a two-phase framework that integrates both paradigms.
In Phase~I, the Fuzzy Best-Worst Method and Fuzzy TOPSIS translate expert judgments across 16~criteria into supplier closeness coefficients.
In Phase~II, these coefficients are embedded in the state representation of a Proximal Policy Optimization agent, creating a prior-augmented policy trained over $3\!\times\!10^{6}$~timesteps.
The framework is evaluated against a vanilla PPO baseline and a deterministic heuristic across three demand and disruption scenarios of increasing severity, totalling 27~independent runs.
Results reveal a regime-dependent trade-off: under stable conditions the prior-augmented agent achieves the lowest procurement cost (Cohen's $d = 1.59$, large effect); under systemic shocks it delivers the highest sustainability score ($d = 0.86$) with markedly lower cross-seed variance; yet under symmetric high-volatility demand, the additional state dimensions impose a learning overhead that favours the simpler baseline.
Both RL agents reduce costs by 54 to 80\,\% relative to the heuristic and maintain bullwhip ratios below unity, actively dampening demand variability.
These findings characterise MCDM-derived priors as a regime-activated asset for RL-based supply chain management.
\end{abstract}

% Add your keywords (separated by \sep)
% JCP allows 1--7 keywords. We select exactly 7, removing the generic
% "Multi-criteria decision-making" (already implied by the two specific
% method names) and "Supply chain resilience" (subsumed by "Sustainable
% supplier selection" and the title). This maximises discoverability
% while respecting the journal's guideline to avoid redundant broad terms.
\begin{keyword}
Sustainable supplier selection \sep Dynamic order allocation \sep Fuzzy Best-Worst Method \sep Fuzzy TOPSIS \sep Proximal Policy Optimization \sep Reinforcement learning \sep Bullwhip effect
\end{keyword}